% =====================================================================
%  Lugworm: Exercises
%  A sequenced, dry-run course for learning Grasshopper and the
%  Lugworm plugin together. No printer required.
%
%  Compile with pdflatex (run twice for the table of contents / links):
%    pdflatex exercises.tex
%    pdflatex exercises.tex
% =====================================================================
\documentclass[11pt]{article}

\usepackage[a4paper,margin=2.2cm]{geometry}
\usepackage[utf8]{inputenc}
\usepackage[T1]{fontenc}
\usepackage{textcomp}
\usepackage{amsmath}
\usepackage{booktabs}
\usepackage{longtable}
\usepackage{array}
\usepackage{xcolor}
\usepackage{enumitem}
\usepackage{parskip}
\usepackage{titlesec}
\usepackage{fancyhdr}
\usepackage{hyperref}

\definecolor{lugwormblue}{RGB}{0,70,140}
\hypersetup{
  colorlinks=true,
  linkcolor=lugwormblue,
  citecolor=lugwormblue,
  urlcolor=lugwormblue,
  pdftitle={Lugworm Exercises},
  pdfauthor={Jasper Vollmert, IDAC, Hamburg University of Technology},
  bookmarksopen=true
}

\titleformat{\section}{\normalfont\Large\bfseries\color{lugwormblue}}{\thesection}{0.6em}{}
\titleformat{\subsection}{\normalfont\large\bfseries}{\thesubsection}{0.6em}{}

\pagestyle{fancy}
\fancyhf{}
\fancyhead[L]{\small Lugworm: Exercises}
\fancyhead[R]{\small\leftmark}
\fancyfoot[C]{\thepage}
\renewcommand{\headrulewidth}{0.4pt}

\newcommand{\param}[1]{\texttt{#1}}
\newcommand{\iamdeg}{$^{\circ}$}

\setlist[itemize]{leftmargin=1.4em, itemsep=0.15em, topsep=0.2em}
\setlist[enumerate]{leftmargin=1.6em, itemsep=0.15em, topsep=0.2em}

% one "exercise" environment: bold run-in labels, tight spacing
\newcommand{\exlabel}[1]{\par\vspace{0.3em}\noindent\textbf{#1:}\enspace}

\title{\color{lugwormblue}\textbf{Lugworm}\\[0.3em]\Large Exercises}
\author{Jasper Vollmert \\ \small Institute of Digital and Autonomous Construction (IDAC) \\ \small Hamburg University of Technology \\ \small jasper.vollmert@tuhh.de}
\date{\small \today}

\begin{document}
\maketitle

\begin{abstract}
\noindent This course teaches Grasshopper and the Lugworm plugin
together, entirely inside Grasshopper. No printer is required for any
exercise. Each session adds one or two components or inputs to what is
already known, using geometry built in Grasshopper or generated with
\param{Test Geometry}. Every result is read from component
outputs: \param{info}, \param{preview}, \param{travels}, \param{layers},
\param{length}, \param{volume}, \param{time}, \param{reachOK}, and
\param{problems}.

The course builds toward one shared final project that combines
everything learned and is printed together once every session is
complete. When the course is taught to a group, each session fills one
lecture or lab block. Companion documents: \emph{beginner-guide.tex}
explains each component in plain language, and \emph{reference.tex}
explains the reasoning and serves as a lookup during class.
\end{abstract}

\tableofcontents
\newpage

% =====================================================================
\section{Structure of each exercise}
% =====================================================================

Every exercise has the same parts: \textbf{Grasshopper concept},
\textbf{Goal}, \textbf{Build}, \textbf{Try this}, \textbf{Read the
output}, and \textbf{Check yourself} (a pass or fail criterion that can
be verified without a printer).

% =====================================================================
\section{Session 0: Wiring the pipeline}
% =====================================================================

\exlabel{Grasshopper concept} finding plugin components on the
\textbf{Lugworm} tab; Panels; Heavy components and the manual compute
button.

\exlabel{Goal} get geometry through \param{Slice}, \param{Lugworm
Place}, and \param{G-code Generator} once, so that every later
exercise only changes numbers.

\exlabel{Build}
\begin{enumerate}
  \item Place an \param{Test Geometry} component (tab \textbf{Lugworm}, panel \textbf{6 Calibrate}). Right-click the component and choose \textbf{Cylinder} from the dropdown. Set \param{height = 80} and \param{radius = 40}.
  \item Place \param{Slice} (\textbf{1 Slice}), \param{Place} (\textbf{2 Place}), and \param{G-code Generator} (\textbf{3 G-code}).
  \item Connect \param{Test Geometry.geometry} to \param{Slice.geo}, \param{Slice.contours} to \param{Place.contours}, and \param{Place.moved} to \param{G-code Generator.crvs}.
  \item Attach \textbf{Panels} to \param{Slice.info}, \param{G-code Generator.info}, and \param{G-code Generator.gcode}. A panel on \param{gcode} shows one line per item. For a large job, prefer the \param{info} summary or save the file, because a panel holding an entire print is slow to draw.
  \item Note that \param{Slice} and \param{G-code Generator} each show a small compute button in the lower-right corner. Both are Heavy components (see \emph{reference.tex} \S 2). Click each button once.
\end{enumerate}

\exlabel{Try this} leave every input at the default for now and click
both compute buttons.

\exlabel{Read the output} \param{Slice.info} should report a
layer count and no per-object errors. \param{Lugworm G-code
Generator.gcode} should show G-code that starts with a header comment
block.

\exlabel{Check yourself} \param{Slice.layers} multiplied by
\param{layerHeight} (default \param{3.0}) should be roughly the
cylinder height of \param{80}, within one layer. If \param{layers} is
\param{0}, check the wiring and remember to click the compute button
after every change. Heavy components do not refresh on their own.

% =====================================================================
\section{Session 1: Layer height, walls, and layer spacing}
% =====================================================================

\exlabel{Grasshopper concept} Number Sliders; clicking the compute
buttons again after a slider change.

\exlabel{Goal} understand how \param{layerHeight}, \param{wallCount},
and \param{spacing} change the sliced result, and connect that to the
physical layer size on \param{G-code Generator}.

\exlabel{Try this}
\begin{enumerate}
  \item Set \param{Slice.layerHeight} to \param{2}, then \param{4}, then \param{8}. Recompute each time and keep \param{G-code Generator.layerHeight} at the same value. Watch \param{layers} change.
  \item Set \param{wallCount} to \param{1}, then \param{3}. Look at \param{preview} in the viewport and count the concentric rings per layer.
  \item Set \param{spacing} to \param{6}, then \param{15}, with \param{wallCount = 3}. Watch how far apart the wall rings sit. The value of \param{spacing} must be greater than 0; otherwise the component raises a hard error.
\end{enumerate}

\exlabel{Read the output} \param{preview} from \param{Lugworm G-code
Generator} shows the rings per layer.

\exlabel{Check yourself} doubling \param{layerHeight} should roughly
halve \param{layers}. With \param{wallCount = 3}, three distinct
concentric curves per layer should be visible, spaced by
\param{spacing}.

% =====================================================================
\section{Session 2: Infill patterns side by side}
% =====================================================================

\exlabel{Grasshopper concept} sweeping an integer slider through named
options; comparing outputs across runs.

\exlabel{Goal} see the practical difference between the four infill
types, and learn to read \param{travels} and \param{time} as an
indicator of print-quality risk.

\exlabel{Try this} on \param{Slice}, set \param{infill} to
\param{1} (zigzag), recompute, and note the extent of \param{Lugworm
G-code Generator.travels} and the value of \param{time}. Repeat for
\param{2} (concentric), \param{3} (spiral), and \param{4} (fermat).

\exlabel{Read the output} compare the four runs. Does \param{time} go
up or down? Do the \param{travels} curves become sparser? Run
\param{Print Stats} (\textbf{5 Analyze}), connected to
\param{G-code Generator.gcode}, and compare the \param{timeMin}
and \param{travelLen} outputs across the four infill choices.

\exlabel{Check yourself} fermat and spiral should show noticeably fewer
and shorter travel moves than zigzag and concentric, because they are
built as continuous space-filling paths rather than as many separate
rings or lines.

% =====================================================================
\section{Session 3: Vase mode (spiralize)}
% =====================================================================

\exlabel{Grasshopper concept} Boolean and menu toggles as mode
switches.

\exlabel{Goal} understand \param{spiralize} as a special case that
requires one closed contour per layer, and see what happens when a
shape does not qualify.

\exlabel{Build} on \param{Test Geometry}, switch the dropdown to
\textbf{Vase} or \textbf{TwistedVase} and set \param{height},
\param{radius}, and the profile inputs (\param{waveAmp},
\param{waveCount}, \param{twistDeg}) as preferred. Set \param{Lugworm
Slice.wallCount = 1} and \param{infill = 0}.

\exlabel{Try this} enable the \textbf{Spiralize} menu toggle on
\param{Slice} and recompute.

\exlabel{Read the output} \param{Slice.info} should confirm that
spiralize ran.

\exlabel{Check yourself} if \param{info} reports instead that spiralize
could not run because a layer does not reduce to one closed contour,
simplify the profile (lower \param{waveAmp} or \param{twistDeg}) until
every horizontal slice is a single closed loop.

% =====================================================================
\section{Session 4: Seam control}
% =====================================================================

\exlabel{Grasshopper concept} a second-stage component (\param{Lugworm
Place}) transforming the tree of a first stage; visualising a point
tree (\param{seams}).

\exlabel{Goal} see how \param{seamMode} changes where the start and
stop point of each layer sits.

\exlabel{Build} reuse the cylinder from Session 1 (planar, several
closed layers).

\exlabel{Try this} cycle \param{Place.seamMode} through
\param{0} (Keep), \param{1} (Align, with a few \param{seamAngle}
values), \param{2} (Stagger, with \param{seamAngle = 15}), \param{3}
(Nearest, with a \param{seamPoint} placed to one side), and,
optionally, \param{4} (Guide, with a \param{seamCurve}) and \param{5}
(Track-below, which keeps the seam aligned with the layer beneath
rather than with a fixed target).

\exlabel{Read the output} display \param{Place.seams} as points
and look from above.

\exlabel{Check yourself} mode \param{1} lines all seam points up in a
column, mode \param{2} rotates them by \param{seamAngle} per layer, and
mode \param{3} clusters them near the \param{seamPoint}. If nothing
changes, confirm that \param{seamAngle}, \param{seamPoint}, or
\param{seamCurve} is being changed, not only \param{seamMode}.

% =====================================================================
\section{Session 5: Flow and pressure basics}
% =====================================================================

\exlabel{Grasshopper concept} reading raw G-code text as a debugging
tool.

\exlabel{Goal} learn where the flow and pressure inputs act by finding
them directly in the \param{gcode} output.

\exlabel{Try this} on a part with \param{wallCount} of 2 or more, so
that both contour ends inside a branch and travel moves are available
for inspection:
\begin{enumerate}
  \item Set \param{pressureOffset = 1} and \param{retractSpeed = 5}. Find a retract and prime pair in the \param{gcode} output and confirm that the pull-back distance and the feed rate match.
  \item Set \param{coast = 1}. Find a contour end before a travel move and confirm that extrusion stops slightly early.
  \item Set \param{leadIn = 3} and \param{leadInBoost = 1.2}. Look for the flow ramp at the start of a printed segment.
  \item Set \param{minLayerTime = 5} and check whether the feed rates of any layer are scaled down as a result.
\end{enumerate}

\exlabel{Read the output} point to the specific lines in \param{gcode}
that changed for each setting.

\exlabel{Check yourself} if a setting seems to change nothing, check
whether the relevant toggle group (\param{advance} or
\param{retractaccel} for firmware hand-off features) is enabled.
Several inputs only take effect once their group is switched on.

% =====================================================================
\section{Session 6: Overhang analysis}
% =====================================================================

\exlabel{Grasshopper concept} reading a warning threshold as a design
signal rather than as an error to dismiss.

\exlabel{Goal} use \param{Overhang} (\textbf{5 Analyze}) to find
where a shape is likely to be unprintable, and understand that the
thresholds are placeholders, not calibrated values.

\exlabel{Build} a shape with a real overhang, for example the
\textbf{Frustum} of \param{Test Geometry} flipped so that the
shape narrows going down, or a \textbf{TwistedVase} with a large
\param{waveAmp}.

\exlabel{Try this} connect the sliced or placed contours to
\param{Overhang} and recompute (the component is Heavy). Look at
the colour-coded preview and the \param{angles} output.

\exlabel{Check yourself} segments beyond \param{warnDeg} (default
25\iamdeg) should be orange, segments beyond \param{failDeg} (default
45\iamdeg) red. As \emph{reference.tex} \S 4.1 explains, both numbers
are uncalibrated placeholders. The exercise is about reading the
analysis, not about trusting the specific angles.

% =====================================================================
\section{Session 7: Print Stats as a pre-flight check}
% =====================================================================

\exlabel{Grasshopper concept} treating a summary output as a gate
before a long operation.

\exlabel{Goal} become comfortable reading \param{Print Stats}
before a file is ever sent to a machine.

\exlabel{Try this} connect \param{G-code Generator.gcode} to
\param{Print Stats}. Set \param{ramDiameter} to the value of
the G-code Generator, so that the volume is computed from the actual E
values rather than from the geometric fallback. Compare
\param{timeMin}, \param{volumeCm3}, and \param{massG} with rough hand
estimates for the size of the part.

\exlabel{Check yourself} deliberately mismatch \param{Lugworm
Slice.layerHeight} and \param{G-code Generator.layerHeight} and
watch \param{volumeCm3} become implausible. The exercise shows why the
two values must match.

% =====================================================================
\section{Session 8: Non-planar modes (advanced)}
% =====================================================================

\exlabel{Grasshopper concept} reading a warning as a design constraint
rather than as a bug to work around.

\exlabel{Part A, conical (\param{mode = 1})} on a shape with a shallow
dome or taper, set \param{mode = 1}, \param{coneAngle = 20}, and
\param{nozzleClear = 25}. If \param{info} warns that the nozzle
clearance is exceeded, lower \param{coneAngle} until the warning
disappears. Do not raise \param{nozzleClear} instead: that hides a real
collision risk rather than resolving the risk (see \emph{reference.tex}
\S 2.1).

\exlabel{Part B, variable layer height} slice the same shape again with
\param{coneAngle = 0} and with the \textbf{Adaptive} group on. Enable
the \param{varh} group on \param{G-code Generator} in both cases
and compare the reported volume with and without it. Neither adaptive
nor conical layers keep a constant vertical spacing, so the extruded
amount is only nominal until \param{varh} is on.

\exlabel{Check yourself} open a definition saved with the retired heat
mode (\param{mode = 2}, removed in 4.0.0) and confirm that
\param{Slice} reports ``heat mode was retired in 4.0.0 - sliced
planar'' instead of failing.

% =====================================================================
\section{Session 9: Robot simulation (dry run)}
% =====================================================================

\exlabel{Grasshopper concept} treating a hard component error as
non-negotiable, unlike a warning.

\exlabel{Goal} learn to validate a print on the robot before the print
runs on real hardware.

\exlabel{Try this}
\begin{enumerate}
  \item Place an \param{Robot Definition} (\textbf{4 Robot}) and select the profile \param{potterbot-scara-v4} (or enter \texttt{"scara"} in a text panel).
  \item Connect the definition and \param{G-code Generator.gcode} to \param{Robot Simulation} (\textbf{4 Robot}). Scrub \param{sim} from \param{0} to \param{1} and watch the simulated arm trace the print. The printed layer and the travel moves are drawn in different colours.
  \item Deliberately place the part (with \param{Place.printCenter}) somewhere unreasonable, far from the bed origin, and recompute.
  \item Switch the definition to \param{none} and note that the component becomes a plain toolpath viewer: no arm, no reach checks, and \param{reachOK} true.
\end{enumerate}

\exlabel{Read the output} watch \param{reachOK}, \param{problems},
\param{messages}, and the message balloon of the component.

\exlabel{Check yourself} an out-of-reach placement should set
\param{reachOK} to false, fill \param{problems}, and raise a red error
on the component (``targets are outside robot limits''). Move the part
back to a sensible position and confirm that all three clear.

% =====================================================================
\section{Session 10: A first pass through calibration (DoE)}
% =====================================================================

\exlabel{Grasshopper concept} treating parameter tuning as a designed
experiment instead of trial and error.

\exlabel{Goal} run one small but complete pass through the calibration
loop without a printer, using invented measurements to see how the
pieces fit together.

\exlabel{Try this}
\begin{enumerate}
  \item \param{DoE Design}: enter two or three factors (for example \param{printSpeed} and \param{pressureOffset}) with reasonable low and high bounds, and choose the design \texttt{"BBD"}. Export the run table.
  \item Invent a plausible quality score for each run (or estimate one from the G-code text of Session 5) and enter the scores into the exported table.
  \item \param{RSM Evaluate}: point \param{tablePath} at the completed table, take \param{factorNames} from the DoE Design output, and set \param{model = "Quadratic"}. Check the reported R\textsuperscript{2} and adjusted R\textsuperscript{2}. With invented data the fit will probably be poor, and that is the point: this is what a weak-fit warning looks like.
  \item \param{Desirability}: connect \param{models}, set a goal string such as \texttt{"max"} for the quality score, and read the recommended settings.
\end{enumerate}

\exlabel{Check yourself} explain in your own words why a real
calibration campaign needs measured prints rather than invented
numbers. This session is about the mechanics, not about a real result.

% =====================================================================
\section{Optional exercises}
% =====================================================================

\begin{itemize}
  \item Compare the \param{arcs} toggle group on and off on a curved part and count the G-code lines before and after.
  \item Turn on \param{skirt} and find the purge ring in \param{preview}.
  \item Turn on \param{floors} with a few base layers and find where the behaviour changes in the \param{gcode} output.
  \item Place an \param{Print Bed}, connect the bed to \param{Place} and \param{Robot Simulation}, and confirm that the part lands on the print centre of the bed without setting \param{printCenter} by hand.
  \item Connect the robot definition to \param{Print Stats} and compare the Cartesian time estimate with the joint-limited estimate near the edge of the SCARA reach.
  \item Use \param{G-code Restart} to cut a print at 50\,\% and inspect the \param{printed} and \param{remaining} curve split.
\end{itemize}

% =====================================================================
\section{Final project (printed together)}
% =====================================================================

Once every session is complete, design one shared part that exercises
at least \textbf{four} of the following, and bring the Grasshopper file
to the final session:

\begin{itemize}
  \item A non-trivial \param{infill} choice, justified by the shape.
  \item A deliberate \param{seamMode} choice, justified by where the seam is least visible.
  \item At least one pressure or flow setting tuned beyond the default (\param{pressureOffset}, \param{coast}, \param{leadIn}).
  \item Adaptive layers or conical mode, justified by a curved or sloped feature, with the matching \param{varh} setting where applicable.
  \item A pass through \param{Overhang} and \param{Print Stats} with no unaddressed red segments and a plausible time estimate.
  \item For printing on the SCARA: a clean \param{Robot Simulation} validation with the matching \param{Robot Definition} connected, \param{reachOK} true, and no red error.
\end{itemize}

Bring the \param{info} panels, the \param{Print Stats} report,
and the robot validation where applicable. The group reviews them
together, and the final files are printed on the shared machine as the
course capstone.

\end{document}